% =========================================================================
% KAPITEL 5.4 -- GLEITOBJEKTE
% 25.08.2026. Elf Objekte: eine Abbildung, sieben Quelltexte, drei Tabellen.
%
% REIHENFOLGE IM TEXT
%   Absatz  2  -> abb:foldstruktur, lst:foldmasken        [lst KUERZBAR]
%   Absatz  4  -> lst:folds
%   Absatz  8  -> lst:merkmalsliste                       [KUERZBAR]
%   Absatz  9  -> lst:dtypes
%   Absatz 11  -> tab:baselines-menge, lst:poisson
%   Absatz 12  -> tab:baselines-struktur, lst:logit
%   Absatz 16  -> lst:test-folds                          [KUERZBAR]
%   ENDE       -> tab:steckbrief
%
% CODEKOMMENTARE SIND NACHGEZOGEN. Die neuen Kommentarzeilen stehen
% seit dem 25.08.2026 auch im Repo (prep/s2_datensaetze.py, prep/config.py,
% vorpruefung/v1_baselines.py); die Ausschnitte unten stimmen mit den Dateien
% ueberein. Die Zeilennummern in den Beschriftungen sind auf den Stand NACH
% dieser Aenderung gesetzt. Liste der Aenderungen am Ende dieser Datei.
%
% Die Abbildung braucht das Paket float (wegen [H]) -- ist in Kapitel 4
% bereits im Einsatz. Die Tabellen brauchen tabularx und booktabs, ebenfalls
% vorhanden.
% =========================================================================


% ---------------------------- AB HIER KOPIEREN ---------------------------

% ---------- Absatz 2 -----------------------------------------------------

\begin{figure}[H]
    \centering
    \includegraphics[width=\textwidth]{Abbildungen/a2_foldstruktur.pdf}
    \caption{Zuteilung der 35 Stadtteile auf fünf Folds und das Hold-out}
    \label{abb:foldstruktur}
\end{figure}


\begin{lstlisting}[style=pythonstil,
  label={lst:foldmasken},
  caption={Trainings- und Testmenge eines Folds
           (\texttt{prep/s2\_datensaetze.py}, Z.~96--111, gekürzt)}]
def fold_masken(daten: pd.DataFrame, k: int) -> tuple[pd.Series, pd.Series]:
    test = daten["fold"] == k
    assert test.any(), (f"Kein Fold {k} im Datensatz "
                        f"(vorhanden: {sorted(daten['fold'].unique())}).")
    train = (daten["fold"] != k) & (daten["ist_holdout"] == 0)
    return train, test
\end{lstlisting}


% ---------- Absatz 4 -----------------------------------------------------

\begin{lstlisting}[style=pythonstil,
  label={lst:folds},
  caption={Fold-Zuteilung mit doppelter Stratifizierung
           (\texttt{prep/s2\_datensaetze.py}, Z.~61--93, gekürzt)}]
def ergaenze_aufteilung(daten: pd.DataFrame, versatz: int = 0,
                        selten: pd.Series | None = None) -> pd.DataFrame:
    # Doppelte Stratifizierung (#30): zuerst nach brand-dominierten Monaten,
    # bei Gleichstand nach Bevoelkerung.
    bev = daten.groupby("stadtteil")[EXPOSURE_ROH].mean()
    if selten is None:
        selten = pd.Series(0, index=bev.index)
    ordnung = (pd.DataFrame({"selten": selten.reindex(bev.index).fillna(0),
                             "bev": bev})
                 .sort_values(["selten", "bev"], ascending=False).index)

    # Reihum auf N_FOLDS + 1 Gruppen; Gruppe 0 ist das Hold-out.
    gruppe = {st: (i + versatz) % (N_FOLDS + 1) for i, st in enumerate(ordnung)}

    d = daten.copy()
    d["fold"] = d["stadtteil"].map(gruppe).astype(int)
    d["ist_holdout"] = (d["fold"] == 0).astype(int)
    return d
\end{lstlisting}


% ---------- Absatz 8 -----------------------------------------------------

\begin{lstlisting}[style=pythonstil,
  label={lst:merkmalsliste},
  caption={Merkmalssatz, einmalig hinterlegt
           (\texttt{prep/config.py}, Z.~106--131, gekürzt)}]
PRAEDIKTOREN = [
    "median_haushaltseinkommen", "armutsquote_pct", "akademikerquote_pct",
    "median_miete", "leerstandsquote_pct", "log_bevoelkerung",
    "log_kriminalitaetsindex",
    "anteil_altbau_vor_1940_pct", "anteil_wohngebaeude_pct",
    "anteil_risikogewerbe_pct",
]

# Rohwert: kein Merkmal, aber Offset des Poisson-GLM.
EXPOSURE_ROH = "gesamtbevoelkerung"

# Saison als sin/cos - der Monat als Zahl gaebe Dezember und Januar Abstand 11.
SAISON = ["monat_sin", "monat_cos"]

# Ein Merkmalssatz fuer alle drei Verfahren. Ohne rohes Jahr und
# Stadtteil-ID: Baeume koennen nicht extrapolieren, Ridge schon - das
# verzerrte den Vergleich.
FEATURE_SETS = {
    "S": PRAEDIKTOREN + SAISON,
}
\end{lstlisting}


% ---------- Absatz 9 -----------------------------------------------------

\begin{lstlisting}[style=pythonstil,
  label={lst:dtypes},
  caption={Vereinheitlichung der Datentypen
           (\texttt{prep/s2\_datensaetze.py}, Z.~151--174, gekürzt)}]
def _setze_datentypen(d: pd.DataFrame, merkmale: list[str]) -> pd.DataFrame:
    # Eine einzige nullable Int64-Spalte genuegt, damit X.to_numpy() ein
    # object-Array liefert: sklearn faengt das still ab, XGBoost nicht (#24).
    d = d.copy()
    for c in merkmale:
        d[c] = pd.to_numeric(d[c], errors="coerce").astype("float64")
    for c in ["jahr", "monat", "jahr_monat", ZIELGROESSE,
              "fold", "ist_holdout", EXPOSURE_ROH] + ANZAHLEN:
        if c in d.columns:
            d[c] = pd.to_numeric(d[c], errors="coerce").astype("int64")
    for c in ("stadtteil", ZIELKLASSE):
        if c in d.columns:
            d[c] = d[c].astype(str)
    return d
\end{lstlisting}


% ---------- Absatz 11 ----------------------------------------------------

\begin{table}[htbp]
    \centering
    \caption{Referenzwerte der Mengenzielgrößen, Mittel über 50 Läufe}
    \label{tab:baselines-menge}
    \footnotesize
    \singlespacing
    \setlength{\tabcolsep}{5pt}
    \begin{tabularx}{\textwidth}{@{}>{\raggedright\arraybackslash}p{3.9cm} c >{\raggedright\arraybackslash}X l l@{}}
        \toprule
        \textbf{Zielgröße} & \textbf{Stufe} & \textbf{Referenz} & \textbf{R\textsuperscript{2}} & \textbf{RMSE} \\
        \midrule
        Einsatzhäufigkeit & 1 & Gesamtmittelwert & $-0{,}744 \pm 0{,}325$ & $69{,}93 \pm 1{,}92$ \\
        Einsatzhäufigkeit & 2 & Poisson-Modell mit Offset & $0{,}542 \pm 0{,}082$ & $33{,}98 \pm 3{,}11$ \\
        \addlinespace
        Einsätze je 1.000 Einwohner & 1 & Gesamtmittelwert & $-1{,}054 \pm 0{,}875$ & $7{,}54 \pm 0{,}13$ \\
        Einsätze je 1.000 Einwohner & 2 & Poisson-Modell mit Offset & $0{,}367 \pm 0{,}261$ & $4{,}08 \pm 0{,}62$ \\
        \bottomrule
    \end{tabularx}
\end{table}


\begin{lstlisting}[style=pythonstil,
  label={lst:poisson},
  caption={Poisson-Modell mit Offset, Stufe 2 der Menge
           (\texttt{vorpruefung/v1\_baselines.py}, Z.~74--108, gekürzt)}]
def poisson_glm(train: pd.DataFrame, test: pd.DataFrame,
                merkmale: list[str] | None = None) -> np.ndarray:
    import statsmodels.api as sm

    spalten = MERKMALE if merkmale is None else list(merkmale)
    X_tr = sm.add_constant(train[spalten].astype(float), has_constant="add")
    X_te = sm.add_constant(test[spalten].astype(float),  has_constant="add")
    y_tr = train[ZIELGROESSE].astype(float)

    # log(Bevoelkerung) als OFFSET, Koeffizient fest auf 1: geschaetzt wird
    # die Rate, hochgerechnet wird erst in der Vorhersage (#13).
    off_tr = np.log(train[EXPOSURE_ROH].astype(float))
    off_te = np.log(test[EXPOSURE_ROH].astype(float))

    # kanonischer log-Link, unpenalisiert - kein freier Hyperparameter
    modell = sm.GLM(y_tr, X_tr, family=sm.families.Poisson(),
                    offset=off_tr).fit()
    return np.asarray(modell.predict(X_te, offset=off_te))
\end{lstlisting}


% ---------- Absatz 12 ----------------------------------------------------

\begin{table}[htbp]
    \centering
    \caption{Referenzwerte der überwiegenden Einsatzart, Mittel über 50 Läufe; für die konstante Vorhersage der Stufe 1 ist AUROC nicht definiert}
    \label{tab:baselines-struktur}
    \footnotesize
    \singlespacing
    \setlength{\tabcolsep}{5pt}
    \begin{tabularx}{\textwidth}{@{}c >{\raggedright\arraybackslash}X l l l@{}}
        \toprule
        \textbf{Stufe} & \textbf{Referenz} & \textbf{Macro-F1} & \textbf{Macro-AUROC} & \textbf{Trefferquote} \\
        \midrule
        1 & Häufigste Klasse & 0,223 & --- & 0,806 \\
        2 & Multinomiales Logit & $0{,}297 \pm 0{,}014$ & 0,705 & 0,584 \\
        \bottomrule
    \end{tabularx}
\end{table}


\begin{lstlisting}[style=pythonstil,
  label={lst:logit},
  caption={Multinomiales Logit, Stufe 2 der Struktur
           (\texttt{vorpruefung/v1\_baselines.py}, Z.~111--139, gekürzt)}]
def logit_glm(train: pd.DataFrame, merkmale: list[str] | None = None):
    from sklearn.linear_model import LogisticRegression
    from sklearn.pipeline import make_pipeline
    from sklearn.preprocessing import StandardScaler

    spalten = MERKMALE if merkmale is None else list(merkmale)

    # unpenalisiert (C = inf); ausgeglichene Klassengewichte statt Resampling:
    # keine duplizierte und keine geloeschte Zeile
    return make_pipeline(
        StandardScaler(),
        LogisticRegression(max_iter=2000, C=np.inf, class_weight="balanced")
    ).fit(train[spalten].astype(float), train[ZIELKLASSE])
\end{lstlisting}


% ---------- Absatz 16 ----------------------------------------------------

\begin{lstlisting}[style=pythonstil,
  label={lst:test-folds},
  caption={Prüfung der Stadtteiltrennung an der erzeugten Datei
           (\texttt{tests/test\_aufbereitung.py}, Z.~134--156, gekürzt)}]
def test_folds_ordnung_und_holdout():
    d = regression()
    holdout = set(d.loc[d["ist_holdout"] == 1, "stadtteil"])
    gesehen = set()
    for k in range(1, N_FOLDS + 1):
        train, test = fold_masken(d, k)
        st_train = set(d.loc[train, "stadtteil"])
        st_test = set(d.loc[test, "stadtteil"])
        assert st_train.isdisjoint(st_test), \
            f"Fold {k}: Stadtteil in Training UND Test: {st_train & st_test}"
        assert st_test.isdisjoint(holdout), f"Fold {k} greift ins Hold-out"
        assert gesehen.isdisjoint(st_test), "Stadtteil in zwei Folds im Test"
        gesehen |= st_test
    assert gesehen | holdout == set(d["stadtteil"]), \
        "Nicht jeder Stadtteil ist genau einmal Testfall oder im Hold-out"
\end{lstlisting}


% ---------- Ende des Abschnitts ------------------------------------------

\begin{table}[htbp]
    \centering
    \caption{Steckbrief der beiden Analysedatensätze}
    \label{tab:steckbrief}
    \footnotesize
    \singlespacing
    \setlength{\tabcolsep}{5pt}
    \begin{tabularx}{\textwidth}{@{}>{\raggedright\arraybackslash}p{4.6cm} >{\raggedright\arraybackslash}X@{}}
        \toprule
        \textbf{Eigenschaft} & \textbf{Ausprägung} \\
        \midrule
        Analyseeinheit & Stadtteil und Monat, in beiden Dateien \\
        Zeitraum & Januar 2015 bis Dezember 2025, 132 Monate \\
        Beobachtungen Menge & 4.620, das sind 35 mal 132 \\
        Beobachtungen Struktur & 4.619, ein Monat ohne Einsatz entfällt \\
        Merkmale & 10 Strukturmerkmale und 2 Saisonterme \\
        Aufteilung & 5 Folds mit 6, 6, 6, 6 und 5 Stadtteilen; 29 Entwicklungs- und 6 Hold-out-Stadtteile \\
        Ausschlüsse & 3 Parkgebiete, 3 Stadtteile ohne durchgängige Zensusabdeckung \\
        Absicherung & 19 automatisierte Prüfungen an den erzeugten Dateien \\
        \bottomrule
    \end{tabularx}
\end{table}

% ---------------------------- BIS HIER KOPIEREN --------------------------


% =========================================================================
% CODEAENDERUNGEN -- am 25.08.2026 im Repo AUSGEFUEHRT, py_compile fehlerfrei
% =========================================================================
% prep/s2_datensaetze.py
%   ergaenze_aufteilung(), vor `bev = ...`
%     # Doppelte Stratifizierung (#30): zuerst nach brand-dominierten Monaten,
%     # bei Gleichstand nach Bevoelkerung.
%   ergaenze_aufteilung(), vor `gruppe = ...` (mit Leerzeile davor)
%     # Reihum auf N_FOLDS + 1 Gruppen; Gruppe 0 ist das Hold-out.
%   _setze_datentypen(), erste Zeile im Rumpf
%     # Eine einzige nullable Int64-Spalte genuegt, damit X.to_numpy() ein
%     # object-Array liefert: sklearn faengt das still ab, XGBoost nicht (#24).
%
% vorpruefung/v1_baselines.py
%   poisson_glm(), vor `off_tr = ...`
%     # log(Bevoelkerung) als OFFSET, Koeffizient fest auf 1: geschaetzt wird
%     # die Rate, hochgerechnet wird erst in der Vorhersage (#13).
%   poisson_glm(), vor `modell = sm.GLM(...)`
%     # kanonischer log-Link, unpenalisiert - kein freier Hyperparameter
%   logit_glm(), vor `return make_pipeline(`
%     # unpenalisiert (C = inf); ausgeglichene Klassengewichte statt Resampling:
%     # keine duplizierte und keine geloeschte Zeile
%
% Zwei alte Punkte gleich miterledigt:
%   prep/config.py Z. 114  "Deskription 5.1" -> "Deskription 4.1"
%   tests/test_aufbereitung.py  "NegBin-Offset" -> "Poisson-Offset"
%
% Durch die acht eingefuegten Zeilen verschieben sich Zeilennummern:
%   fold_masken        92-107  ->  96-111
%   ergaenze_aufteilung 61-90  ->  61-93
%   _setze_datentypen 147-170  -> 151-174
%   poisson_glm         74-105 ->  74-108
%   logit_glm          107-132 -> 111-139
% Die Beschriftungen oben tragen bereits die neuen Nummern. Wenn du im Repo
% noch etwas aenderst, hier nachziehen.
%
% NOCH OFFEN im Code (nicht 5.4-relevant, aber auf der Liste):
%   tests/test_aufbereitung.py Z. 5-6   Verweis auf einen Anhang entfernen
%   prep/config.py Z. 76-77             0-60 min als gesetzte Annahme kennzeichnen
%   prep/s2_datensaetze.py Z. ~380      "Zaehlgroesse" -> "stetigen Groesse"
%   prep/config.py Z. 41                Census-API-Schluessel in eine
%                                       Umgebungsvariable, alten Schluessel
%                                       vor der Abgabe zuruecknehmen
% =========================================================================
